% AgroVoltaic --- documentacion de arquitectura.
% Compilar: pdflatex -interaction=nonstopmode arquitectura.tex (dos veces, por el indice).
\documentclass[11pt,a4paper]{article}
% Preambulo compartido: paquetes, paleta, estilos de TikZ y macros.
% Responsabilidad unica: configurar el documento. No contiene contenido.

% Tipografia (requiere xelatex; ver README). Pagella para el cuerpo, Heros para
% titulos y diagramas, JetBrains Mono para codigo. Las tres estan instaladas:
% las dos primeras vienen con texlive-fontsrecommended, la tercera es del sistema.
\usepackage{fontspec}
\usepackage{unicode-math}
% Las TeX Gyre viven en texmf y NO estan registradas en fontconfig: se cargan por
% nombre de archivo, no por nombre de familia.
\setmainfont{texgyrepagella}[
  Extension=.otf, UprightFont=*-regular, ItalicFont=*-italic,
  BoldFont=*-bold, BoldItalicFont=*-bolditalic]
\setsansfont{texgyreheros}[
  Extension=.otf, UprightFont=*-regular, ItalicFont=*-italic,
  BoldFont=*-bold, BoldItalicFont=*-bolditalic, Scale=MatchLowercase]
\setmonofont{JetBrains Mono}[Scale=MatchLowercase]
\setmathfont{texgyrepagella-math.otf}
\usepackage[spanish,es-noquoting,es-tabla]{babel}
% expansion no existe en xetex; protrusion si, y es la que mas se nota.
\usepackage[protrusion=true,expansion=false]{microtype}
\usepackage{geometry}
\geometry{a4paper,margin=2.4cm,headsep=12pt}
\usepackage{amsmath}
\usepackage{booktabs}
\usepackage{longtable}
\usepackage{array}
\usepackage{enumitem}
\usepackage{tcolorbox}
\usepackage{tikz}
\usepackage{fancyhdr}
\usepackage{titlesec}
\usepackage[hidelinks]{hyperref}

\usetikzlibrary{positioning,shapes.geometric,arrows.meta,calc,fit,backgrounds}
\tcbuselibrary{skins,breakable}

% --- Paleta -----------------------------------------------------------------
\definecolor{solar}{HTML}{B26B00}   % naranja apagado: fuentes fisicas / energia
\definecolor{tierra}{HTML}{4F6B3A}  % verde: agro / AgroDash
\definecolor{agua}{HTML}{20566E}    % azul: almacenamiento
\definecolor{tinta}{HTML}{1F2328}   % texto
\definecolor{humo}{HTML}{6B7280}    % texto secundario
\definecolor{papel}{HTML}{F5F3EE}   % fondo de cajas
\definecolor{borde}{HTML}{D6D2C8}

\hypersetup{colorlinks=false,pdftitle={AgroVoltaic --- Arquitectura},
            pdfauthor={Proyecto AgroVoltaic (TEC)}}

% --- Titulos ----------------------------------------------------------------
\titleformat{\section}{\Large\bfseries\color{agua}}{\thesection}{0.7em}{}
\titleformat{\subsection}{\large\bfseries\color{tinta}}{\thesubsection}{0.7em}{}
\titlespacing*{\section}{0pt}{2.2ex plus 1ex minus .2ex}{1.2ex plus .2ex}

\pagestyle{fancy}
\fancyhf{}
\fancyhead[L]{\small\color{humo}AgroVoltaic --- Arquitectura del sistema}
\fancyhead[R]{\small\color{humo}\thepage}
\renewcommand{\headrulewidth}{0.4pt}

% --- Macros -----------------------------------------------------------------
\newcommand{\cd}[1]{\texttt{\small #1}}          % identificador / ruta / comando
% \nota: texto secundario. OJO: no admite \\ adentro — el salto de linea de un
% nodo de TikZ es una celda de alineacion y no cruza el grupo. Una \nota por linea.
\newcommand{\nota}[1]{{\color{humo}\small #1}}

\newtcolorbox{clave}[1][]{
  enhanced, breakable, colback=papel, colframe=agua, boxrule=0.6pt,
  left=8pt, right=8pt, top=6pt, bottom=6pt, arc=2pt,
  fonttitle=\bfseries\small, coltitle=white, colbacktitle=agua, #1}

\newtcolorbox{aviso}[1][]{
  enhanced, breakable, colback=white, colframe=solar, boxrule=0.6pt,
  left=8pt, right=8pt, top=6pt, bottom=6pt, arc=2pt,
  fonttitle=\bfseries\small, coltitle=white, colbacktitle=solar, #1}

% --- Estilos de diagrama ----------------------------------------------------
\tikzset{
  comp/.style={
    draw=agua, fill=papel, line width=0.7pt, rounded corners=3pt,
    align=center, inner sep=6pt, minimum height=10mm, font=\small},
  compx/.style={comp, draw=humo, fill=white},          % componente externo
  fuente/.style={comp, draw=solar, fill=solar!8},      % fuente fisica / dato crudo
  agro/.style={comp, draw=tierra, fill=tierra!8},      % region Cartago
  almacen/.style={
    draw=agua, fill=agua!8, line width=0.7pt, cylinder, shape border rotate=90,
    aspect=0.16, align=center, inner sep=6pt, minimum height=13mm, font=\small},
  etapa/.style={comp, minimum width=27mm},
  flecha/.style={-{Stealth[length=2.4mm]}, line width=0.7pt, draw=tinta},
  flechad/.style={flecha, dashed, draw=humo},
  etiq/.style={font=\scriptsize\itshape, color=humo, inner sep=2pt, fill=white},
  grupo/.style={draw=borde, line width=0.6pt, rounded corners=4pt, inner sep=8pt},
  % Sin guiones dentro de los diagramas: partir "ana-lizador" en una caja de
  % 27 mm se lee peor que dejar la caja un poco mas holgada.
  every picture/.append style={execute at begin picture={%
    \hyphenpenalty=10000\relax\exhyphenpenalty=10000\relax}},
  titgrupo/.style={font=\scriptsize\bfseries, color=humo},
}


\begin{document}

\begin{titlepage}
\centering
\vspace*{3.5cm}
{\color{humo}\large PROYECTO AGROVOLTAICO --- TEC COSTA RICA\par}
\vspace{1.2cm}
{\Huge\bfseries\color{agua} AgroVoltaic\par}
\vspace{0.4cm}
{\LARGE Arquitectura del sistema\par}
\vspace{1.2cm}
{\large\color{tinta} Del sensor en el campo al agente que responde preguntas\par}
\vfill
\begin{minipage}{0.82\textwidth}
\small\color{tinta}
Plataforma de datos del sistema agrovoltaico de San Carlos: el ETL que estandariza
diecinueve meses de CSV crudos, dos agentes con modelo de lenguaje ---uno de análisis
histórico y otro de pronóstico--- y la consola con la que se depuran. Este documento
describe cómo encajan las piezas, qué decisión hay detrás de cada una y qué queda
pendiente.
\end{minipage}
\vfill
{\color{humo} \today\par}
\end{titlepage}

\tableofcontents
\newpage

\section{Panorama}

AgroVoltaic es la plataforma de datos del sistema agrovoltaico del TEC en San Carlos
(Costa Rica): paneles solares bifaciales instalados sobre cultivo. El repositorio es un
\emph{monorepo} con cuatro piezas que se pueden levantar por separado.

\begin{center}
\small
\begin{tabular}{@{}p{34mm}p{34mm}p{72mm}@{}}
\toprule
\textbf{Pieza} & \textbf{Dónde vive} & \textbf{Qué resuelve} \\
\midrule
ETL de CSV & \cd{src/agrovoltaic} & Estandarizar 19 meses de CSV crudos del inversor
y los sensores, y cargarlos a Supabase de forma idempotente. \\
Agente de pronóstico & \cd{agente-pronostico/} & Pronosticar irradiancia y humedad de
suelo con anclaje físico; el LLM sólo orquesta. \\
Agente analizador & \cd{agente-analizador/} & Responder preguntas sobre el histórico
fotovoltaico ya limpio, componiendo herramientas atómicas. \\
Consola de depuración & \cd{mvp-debugger/} & Ver la traza completa de cada agente y
cruzar cada número contra los datos. \\
\bottomrule
\end{tabular}
\end{center}

\subsection{El problema de origen}

Los datos no tenían estándar. El EDA sobre los 285 CSV encontró \textbf{13 esquemas
distintos}: nombres de columna que cambian entre cortos (\cd{vpv1}), largos
(\cd{Voltaje PV1 [V]}) y \cd{snake\_case}; erratas persistentes (\cd{Energì} con acento
grave en 72 archivos, \cd{POTencia}, \cd{Corriente PV2[A]} sin espacio); filas de
sensores distintos intercaladas en el mismo archivo; irradiancia sin calibrar con un
\emph{offset} nocturno constante de $-38{,}845$; temperaturas saturadas en 85\,°C
(el código de error clásico del DS18B20 desconectado); intervalos de muestreo que van
de 2 segundos a 5 minutos según la época; y dos huecos largos de 126 y 71 días.

\begin{clave}[title=Regla rectora del tratamiento de datos]
Validada con Leo Cardinale el 2026-08-10: \textbf{el dato crudo se guarda tal cual en la
base; toda corrección vive en una capa de análisis} (vistas SQL) que genera variables
nuevas. Ninguna transformación destructiva en la ingesta. Esta regla reemplazó las
decisiones previas de convertir 85\,°C a NULL, forzar el \emph{offset} a cero y
remuestrear todo a 5 minutos durante la carga.
\end{clave}

\subsection{Dos regiones, sin base central}

El proyecto abarca dos sitios y \textbf{no} se fusionan sus bases:

\begin{itemize}[leftmargin=*,itemsep=2pt]
  \item \textbf{Cartago} --- recolección automática hacia \textbf{AgroDash} (PostgreSQL,
    app Rust/Axum). Sensores de suelo y ambiente: humedad, conductividad, temperatura,
    irradiancia, PAR. No mide fotovoltaica.
  \item \textbf{San Carlos} --- recolección semi-manual de CSV hacia \textbf{Supabase}.
    Es la región fotovoltaica.
\end{itemize}

El matiz que define la arquitectura: \textbf{San Carlos está partido entre las dos
bases}. Lo eléctrico (voltaje, corriente, potencia, energía, irradiancia del inversor)
vive en Supabase; lo ambiental y de suelo de San Carlos vive en AgroDash, en cajas con
sufijo \cd{SC}. Como Cartago no mide fotovoltaica, no existe comparación cruzada
PV contra PV: lo único común entre regiones son las variables ambientales.

% Pagina propia, sin rotar: el diagrama esta dibujado en vertical justamente para
% que entre derecho y se lea sin ladear la cabeza.
\begin{figure}[p]
\centering
\resizebox{\textwidth}{!}{% Figura: arquitectura global. Flujo de arriba hacia abajo, una banda por etapa:
% fuentes -> acarreo -> ingesta -> almacen -> agentes -> clientes. Vertical a
% proposito: asi entra en una pagina normal sin rotarla ni encogerla hasta
% volverla ilegible.
\begin{tikzpicture}[
  banda/.style={font=\scriptsize\bfseries, color=humo, anchor=east},
  fila/.style={text width=27mm, minimum height=11mm},
]

% --- Coordenadas: una y por banda, una x por columna -------------------------
\def\yfuentes{0}
\def\yacarreo{-3.2}
\def\yingesta{-6.4}
\def\yalmacen{-9.6}
\def\yagentes{-12.8}
\def\yclientes{-16.0}
\def\xpv{3.2}
\def\xamb{10.6}

% --- Banda 1: fuentes fisicas ------------------------------------------------
\node[fuente,fila] (inv) at (0,\yfuentes) {Inversor solar\\\nota{PV1 + PV2}};
\node[fuente,fila] (pir) at (\xpv,\yfuentes) {Piranómetros\\\nota{incidente / reflejada}};
\node[fuente,fila] (ds)  at (6.4,\yfuentes) {DS18B20\\\nota{temp.\ de arreglo}};
\node[agro,fila] (cajas) at (\xamb,\yfuentes) {Cajas y nodos SC\\\nota{suelo, ambiente}};

\node[grupo, fit=(inv)(pir)(ds),
      label={[titgrupo]above:SAN CARLOS --- FOTOVOLTAICO}] (gf) {};
\node[grupo, fit=(cajas),
      label={[titgrupo]above:SAN CARLOS --- AMBIENTAL}] (gc) {};

% --- Banda 2: acarreo --------------------------------------------------------
\node[comp,fila] (csv) at (\xpv,\yacarreo)
  {CSV crudos\\\nota{285 archivos}\\\nota{13 esquemas}};
\node[agro,fila] (agrodash) at (\xamb,\yacarreo)
  {AgroDash\\\nota{PostgreSQL}\\\nota{región Cartago}};

% --- Banda 3: ingesta --------------------------------------------------------
\node[comp,fila] (etlcsv) at (\xpv,\yingesta)
  {\textbf{ETL de CSV}\\\cd{src/agrovoltaic}\\\nota{lote, manual}};
\node[comp,fila] (etlamb) at (\xamb,\yingesta)
  {\textbf{ETL ambiental}\\\cd{pronostico.etl}\\\nota{timer, 15 min}};

% --- Banda 4: almacenamiento -------------------------------------------------
\node[almacen,text width=46mm] (store) at (6.9,\yalmacen)
  {\textbf{Supabase AgroVoltaic}\\\nota{crudo + vistas + store del agente}};

% --- Banda 5: agentes --------------------------------------------------------
\node[comp,fila] (ana) at (\xpv,\yagentes)
  {\textbf{Agente analizador}\\\nota{Q\&A histórico PV}\\\nota{\cd{:8010}}};
\node[comp,fila] (pro) at (\xamb,\yagentes)
  {\textbf{Agente de pronóstico}\\\nota{irradiancia + humedad}\\\nota{\cd{:8000}}};
\node[compx,text width=22mm] (llm) at (6.9,\yagentes) {API de Claude};

% --- Banda 6: clientes -------------------------------------------------------
\node[comp,fila] (web) at (4.6,\yclientes)
  {\textbf{Consola}\\\cd{mvp-debugger}\\\nota{Next.js, local}};
\node[compx,fila] (vf) at (9.2,\yclientes) {VisioneFlow\\\nota{flujos programados}};

% --- Etiquetas de banda ------------------------------------------------------
\node[banda] at (-1.9,\yfuentes)  {FUENTES};
\node[banda] at (-1.9,\yacarreo)  {ACARREO};
\node[banda] at (-1.9,\yingesta)  {INGESTA};
\node[banda] at (-1.9,\yalmacen)  {ALMACÉN};
\node[banda] at (-1.9,\yagentes)  {AGENTES};
\node[banda] at (-1.9,\yclientes) {CLIENTES};

% --- Flujos ------------------------------------------------------------------
\draw[flecha] (inv.south) -- (csv.north);
\draw[flecha] (pir.south) -- (csv.north);
\draw[flecha] (ds.south)  -- (csv.north);
\draw[flecha] (cajas)     -- (agrodash);

\draw[flecha] (csv)      -- (etlcsv);
\draw[flecha] (agrodash) -- node[etiq,right]{sólo lectura} (etlamb);

\draw[flecha] (etlcsv.south) -- ++(0,-0.6) -| (store.north);
\draw[flecha] (etlamb.south) -- ++(0,-0.6) -| (store.north);

\draw[flecha] (store.south) -- ++(0,-0.6) -| (ana.north);
\draw[flecha] (store.south) -- ++(0,-0.6) -| (pro.north);
\node[etiq,anchor=south west] at ($(ana.north)+(0.12,0.02)$) {sólo lectura};
\node[etiq,anchor=south east] at ($(pro.north)+(-0.12,0.02)$) {lee y escribe predicciones};

\draw[flechad] (ana) -- (llm);
\draw[flechad] (pro) -- (llm);

\draw[flecha]  (ana.south) -- ++(0,-0.6) -| (web.north);
\draw[flecha]  (pro.south) -- ++(0,-0.6) -| (web.north);
\draw[flechad] (pro.south) -- ++(0,-0.6) -| (vf.north);

\end{tikzpicture}
}
\caption{Arquitectura global. Dos rutas de ingesta independientes convergen en una sola
Supabase, que es lo único que los agentes leen.}
\end{figure}

\section{ETL de CSV a Supabase}

Es el pipeline histórico: convierte la carpeta de CSV crudos en dos tablas normalizadas.
No es un script de una sola vez, sino un proceso reproducible que se vuelve a correr
cada vez que aparecen archivos nuevos.

\begin{figure}[htbp]
\centering
\resizebox{\textwidth}{!}{% Figura: pipeline ETL de CSV, etapa por etapa, con el modulo responsable.
\begin{tikzpicture}[node distance=9mm]

\node[fuente,etapa] (a) {CSV crudo\\\nota{ragged, 13 esquemas}};
\node[etapa,right=of a] (b) {\textbf{extract}\\\nota{leer + tipar}\\\nota{+ normalizar cabeceras}};
\node[etapa,right=of b] (c) {\textbf{transform}\\\nota{separar flujos}\\\nota{+ remuestrear}};
\node[etapa,right=of c] (d) {\textbf{load}\\\nota{UPSERT por}\\\nota{\cd{timestamp}}};

\node[almacen,font=\scriptsize,text width=38mm,right=12mm of d] (t1) {\cd{monitoreo\_sc\_electrico}\\\nota{5 min}};
\node[almacen,font=\scriptsize,text width=38mm,below=8mm of t1] (t2) {\cd{radiacion\_sc\_15s}\\\nota{15 s}};

\draw[flecha] (a) -- (b);
\draw[flecha] (b) -- node[etiq,above]{\cd{slugify}} (c);
\draw[flecha] (c) -- (d);
\draw[flecha] (d.east) -- ++(4mm,0) |- (t1);
\draw[flecha] (d.east) -- ++(4mm,0) |- (t2);

% Estado / idempotencia
\node[comp,text width=30mm,below=14mm of b] (log) {\cd{\_ingest\_log}\\\nota{md5 por archivo}};
\draw[flechad] (a) |- node[etiq,pos=0.72,below]{¿cambió?} (log);
\draw[flechad] (log) -| (d.south);

% Capa solar
\node[comp,text width=36mm,below=14mm of t2] (solar) {\textbf{calibracion} + \textbf{performance}\\\nota{pvlib: clear-sky y POA}};
\draw[flecha] (t2) -- (solar);
\node[almacen,font=\scriptsize,text width=36mm,left=12mm of solar] (t3) {\cd{radiacion\_sc\_clearsky}\\ \cd{radiacion\_sc\_poa}};
\draw[flecha] (solar) -- (t3);

\end{tikzpicture}
}
\caption{Pipeline de CSV. El \cd{\_ingest\_log} decide qué archivos ni siquiera se abren;
la capa solar (pvlib) se refresca al final de cada corrida.}
\end{figure}

\subsection{Una sola lista, y todo lo demás se deduce de ella}

El problema de fondo de estos CSV es que el mismo dato viene escrito de muchas maneras
distintas. El voltaje del primer arreglo aparece como \cd{vpv1} en unos archivos, como
\cd{Voltaje PV1 [V]} en otros y como \cd{voltaje\_pv1\_v} en los más nuevos, y a eso hay
que sumarle las erratas y los acentos que se fueron colando. En total, unas setenta formas
distintas de nombrar apenas una docena de cosas.

La salida fácil sería escribir a mano la lista de nombres buenos y repetirla en cada lugar
donde hace falta: al crear las tablas, al insertar, al agrupar por intervalos de tiempo.
Es también la forma más segura de que dentro de seis meses alguien agregue una variable en
tres de esos cuatro lugares, se olvide del cuarto, y nadie se entere hasta que los números
no cuadren.

Por eso acá hay \textbf{un solo lugar} donde se declara qué se mide: un diccionario con
una entrada por cada cosa medida, que dice cuál es su nombre bueno. Todo lo demás sale de
ahí, incluido el SQL que crea la base de datos, que es un archivo \textbf{generado} y no
se escribe a mano.

Para quien tenga que trabajar con esto mañana, la diferencia se nota en dos cosas:

\begin{itemize}[leftmargin=*,itemsep=2pt]
  \item \textbf{Agregar una variable es escribir un renglón.} Aparece sola en el esquema,
    en la tabla y al agrupar por intervalos. No hay que acordarse de nada más.
  \item \textbf{Un nombre nuevo para algo que ya existe no da trabajo.} Si el archivo del
    mes que viene escribe la misma medición con otra errata, otro acento u otras unidades,
    el sistema la reconoce y la manda al mismo sitio, sin tocar una línea de código. Así
    fue como esas setenta variantes quedaron reducidas a una lista limpia.
\end{itemize}

En el código, el diccionario es \cd{CONCEPT\_MAP} y vive en el módulo \cd{normalize}.

\subsection{Módulos}

\begin{center}
\small
\begin{tabular}{@{}p{26mm}p{114mm}@{}}
\toprule
\textbf{Módulo} & \textbf{Responsabilidad única} \\
\midrule
\cd{normalize} & \cd{slugify} normaliza acentos, unidades y mayúsculas; \cd{CONCEPT\_MAP}
  traduce el slug al nombre canónico. \\
\cd{schemas} & Deriva \cd{CANONICAL\_COLUMNS}, infiere etiquetas del nombre y decide el
  método de agregación por etiqueta. \\
\cd{extract} & Lee el CSV tolerando filas irregulares, normaliza cabeceras, tipa y parsea
  el \cd{timestamp}. \\
\cd{transform} & \cd{split\_streams} separa el flujo eléctrico del de radiación
  \emph{por columnas}; remuestrea a 5\,min y 15\,s. \textbf{Sin limpieza}. \\
\cd{ddl} & Genera tablas, vistas de corrección y calibración, diccionario y sentencias RLS. \\
\cd{clearsky} & GHI de cielo despejado con pvlib (modelo Ineichen). \\
\cd{calibracion} & Puebla \cd{radiacion\_sc\_clearsky} de forma incremental. \\
\cd{performance} & Calcula POA por arreglo (transposición Erbs + isotrópica, modelo
  bifacial de dos planos) y puebla \cd{radiacion\_sc\_poa}. \\
\cd{load} & UPSERT por \cd{timestamp} en cada tabla. \\
\cd{state} & Registro md5 en \cd{\_ingest\_log}. \\
\cd{pipeline} & Orquesta \emph{extract $\to$ transform $\to$ load} y el refresco solar. \\
\cd{cli} & Menú interactivo numerado. \\
\bottomrule
\end{tabular}
\end{center}

\subsection{Idempotencia}

Dos niveles, deliberadamente redundantes:

\begin{itemize}[leftmargin=*,itemsep=2pt]
  \item \textbf{Nivel archivo}: \cd{\_ingest\_log} guarda el md5 de cada CSV. Si no
    cambió, ni se abre.
  \item \textbf{Nivel fila}: \cd{timestamp} es PRIMARY KEY con
    \cd{ON CONFLICT DO UPDATE}. Reprocesar nunca duplica, aunque el md5 falle.
\end{itemize}

Un archivo que revienta no frena la corrida: se registra el fallo, se hace \emph{rollback}
de esa transacción y el pipeline sigue con el siguiente. La capa solar va en su propio
\cd{try}, para que un problema de pvlib no invalide una carga de datos que ya funcionó.

\subsection{Uso}

Un solo comando, menú interactivo:

\begin{center}\cd{python3 main.py}\end{center}

Las opciones van en orden de confianza creciente: auditar el dataset, ejecutar en seco
sin tocar la base, generar el DDL, aplicarlo en Supabase, y recién ahí cargar datos
(incremental o reprocesando todo). Para archivos nuevos basta soltarlos en
\cd{dataset/Monitoreo-AgroVoltaic-SC-NEW/} y repetir la carga incremental.

\begin{aviso}[title=Lo que este ETL todavía no hace]
La \textbf{separación fina de filas mezcladas} (el llamado Paso 2) sigue pendiente. Hoy
se conservan las filas que coinciden con la cabecera del archivo y las irregulares se
saltan con registro en el log. El criterio acordado es recuperar lo que se pueda y
aceptar huecos; existe un par de archivos original/corregido como referencia.
\end{aviso}

\section{Modelo de datos}

Todo vive en un único proyecto Supabase (\cd{jijklguopafevyucogro}). El modelo tiene
cuatro capas: el crudo se guarda intacto y cada capa siguiente es una vista que agrega
información sin destruir la anterior.

\begin{figure}[htbp]
\centering
\resizebox{0.98\textwidth}{!}{% Figura: capas del modelo de datos. El crudo nunca se toca; cada capa agrega.
\begin{tikzpicture}[node distance=7mm]

\node[almacen,text width=46mm] (crudo) {\textbf{Capa 0 --- crudo}\\
  \cd{monitoreo\_sc\_electrico}\\ \cd{radiacion\_sc\_15s}\\
  \nota{tal como salió del CSV}};

\node[comp,text width=46mm,right=14mm of crudo] (corr) {\textbf{Capa 1 --- corrección}\\
  \cd{v\_sc\_electrico\_corregido}\\ \cd{v\_sc\_radiacion\_corregida}\\
  \nota{85\,°C, rangos, offset, pre-2025}};

\node[comp,text width=46mm,right=14mm of corr] (cal) {\textbf{Capa 2 --- calibración}\\
  \cd{v\_sc\_radiacion\_calibrada}\\
  \nota{W/m$^2$, kt*, \cd{qc\_ok}}};

\node[comp,text width=46mm,below=12mm of cal] (perf) {\textbf{Capa 3 --- desempeño}\\
  \cd{v\_sc\_performance}\\
  \nota{\cd{pr\_pv1}, \cd{pr\_pv2}}};

\node[comp,text width=42mm,below=12mm of corr] (cs) {\cd{radiacion\_sc\_clearsky}\\ \cd{radiacion\_sc\_poa}\\
  \nota{calculadas con pvlib}};

\draw[flecha] (crudo) -- node[etiq,above]{vista} (corr);
\draw[flecha] (corr) -- node[etiq,above]{vista} (cal);
\draw[flecha] (cs) -- (cal);
\draw[flecha] (cs) -- (perf);
\draw[flecha] (corr.south) |- (perf.west);

\node[etiq,below=3mm of crudo,align=center,fill=none]
  {sin pérdida de información:\\ nada se sobrescribe};

\end{tikzpicture}
}
\caption{Capas del modelo. Una consulta puede pedir el dato tal como salió del sensor o
el dato corregido y calibrado, según lo que necesite justificar.}
\end{figure}

\subsection{Relaciones}

\begin{center}
\small
\begin{tabular}{@{}p{44mm}p{18mm}p{74mm}@{}}
\toprule
\textbf{Relación} & \textbf{Tipo} & \textbf{Contenido} \\
\midrule
\cd{monitoreo\_sc\_electrico} & tabla & Crudo eléctrico a 5\,min: voltaje, corriente,
  potencia y energía por arreglo, salida AC, temperatura del inversor y de cada arreglo. \\
\cd{radiacion\_sc\_15s} & tabla & Crudo de radiación a 15\,s: incidente, reflejada,
  albedo, y las mismas del piranómetro SP722. \\
\cd{radiacion\_sc\_clearsky} & tabla & \cd{cs\_ghi\_wm2} por instante, calculado con pvlib. \\
\cd{radiacion\_sc\_poa} & tabla & Irradiancia en el plano del arreglo, frontal y
  bifacial, para PV1 y PV2. \\
\cd{diccionario\_variables} & tabla & Definición en prosa de cada variable. \\
\cd{\_ingest\_log} & tabla & md5 y conteo por archivo procesado. \\
\midrule
\cd{v\_sc\_electrico\_corregido} & vista & Aplica rangos válidos; fuera de rango pasa a NULL. \\
\cd{v\_sc\_radiacion\_corregida} & vista & Neutraliza el \emph{offset}, recorta negativos
  y anula el período inválido; marca \cd{valido}. \\
\cd{v\_sc\_radiacion\_calibrada} & vista & Radiación en W/m$^2$, \cd{kt\_star} y
  bandera \cd{qc\_ok}. \\
\cd{v\_sc\_performance} & vista & \cd{pr\_pv1} y \cd{pr\_pv2}: Performance Ratio por arreglo. \\
\midrule
\cd{lecturas\_ambientales\_sc} & tabla & Store de la ingesta ambiental, formato largo. \\
\cd{predicciones} & tabla & Auditoría de cada pronóstico emitido. \\
\cd{agente\_log} & tabla & Log estructurado del ETL y de los agentes. \\
\cd{gasto\_diario} & tabla & Gasto acumulado en el LLM, por día. \\
\bottomrule
\end{tabular}
\end{center}

\subsection{Reglas de corrección}

Los umbrales son parámetros en \cd{config.py} y el SQL de las vistas se genera a partir
de ellos, así que cambiar un límite es cambiar una constante y regenerar el DDL.

\begin{center}
\small
\begin{tabular}{@{}p{40mm}p{34mm}p{62mm}@{}}
\toprule
\textbf{Variable} & \textbf{Rango válido} & \textbf{Motivo} \\
\midrule
Temperaturas & 10--80\,°C & Criterio de Leo Cardinale; además el valor exacto 85{,}0
  es el código de error del DS18B20 desconectado. \\
Potencia por string & 0--5000\,W & Cota física provisional; el instalado real es
  1420\,Wp por arreglo. \\
Voltaje de string & 0--600\,V & Rango del inversor. \\
Corriente de string & 0--20\,A & Rango del inversor. \\
Frecuencia de red & 55--65\,Hz & Red de 60\,Hz. \\
Tensión AC & 100--280\,V & Salida monofásica. \\
Albedo & 0--1 & Definición del cociente. \\
Irradiancia & no aplica & \emph{Offset} $-38{,}845 \to 0$; negativos $\to 0$; todo lo anterior al
  2025-07-01 pasa a NULL (error de sensor corregido a mediados de 2025). \\
\bottomrule
\end{tabular}
\end{center}

\subsection{Calibración y desempeño}

No hay constante de calibración de fábrica: lo que se creía una ``celda calibrada'' resultó
ser un nombre comercial. Se calibra entonces contra el cielo despejado, comparando lo que
mide el sensor con lo que debería medir un día sin nubes. Los datos del período válido ya
venían en W/m$^2$, así que la escala quedó en 1{,}0 y lo que aporta esta capa es el índice
de claridad y el control de calidad.

Toda la capa se reduce a dos fórmulas. La primera dice cuánta de la luz posible llegó
de verdad:

\begin{equation*}
  \mathrm{kt^*} = \frac{\mathrm{GHI_{medida}}}{\mathrm{GHI_{cielo\ despejado}}}
  \qquad\text{sólo donde}\quad \mathrm{GHI_{cs}} > 20\ \mathrm{W/m^2}
\end{equation*}

La segunda compara la potencia real con la que cabría esperar de la luz que recibe el
plano del arreglo:

\begin{equation*}
  \mathrm{PR} = \frac{P/P_0}{\mathrm{POA}/1000}
  \qquad P_0 = 1420\ \mathrm{Wp\ por\ arreglo}
\end{equation*}

Como los paneles son bifaciales, la POA suma las dos caras: frontal $+\ \varphi\cdot$trasera,
con $\varphi = 0{,}80$ y el albedo medido.

El resultado da confianza porque converge: dos arreglos con geometrías muy distintas
aterrizan casi en el mismo número, PV1 inclinado en $0{,}622$ y PV2 vertical en $0{,}626$.

\begin{clave}[title=Geometría confirmada del sitio]
1420\,Wp por arreglo ($4 \times 355$\,Wp), 2840\,Wp en total, ambos bifaciales.
PV1 es el arreglo inclinado (\emph{tilt} 20°, azimut 150°); PV2 es el vertical
(\emph{tilt} 90°, azimut 50°), con Norte $=0$° y sentido horario positivo. Zona horaria
UTC$-6$ sin horario de verano.
\end{clave}

\input{secciones/04-pronostico}
\section{Agente analizador}

Responde preguntas en español sobre el histórico fotovoltaico ya limpio. Mismo principio
que el forecaster: el LLM entiende y redacta, los números salen de consultas SQL de sólo
lectura contra las vistas corregidas y calibradas.

\subsection{Una herramienta, un archivo, una pregunta}

La diferencia con el pronóstico está en la granularidad. Aquí cada herramienta responde
\emph{una} pregunta concreta y el LLM \textbf{compone} varias cuando la pregunta es
compuesta. No hay una herramienta genérica de ``ejecutá este SQL'', que sería la vía
rápida y también la que abre la puerta a inyección y a respuestas incomprobables.

\begin{center}
\small
\begin{tabular}{@{}p{26mm}p{114mm}@{}}
\toprule
\textbf{Herramienta} & \textbf{Qué responde} \\
\midrule
\cd{energia} & Energía generada por arreglo (integral de potencia). \\
\cd{performance} & Performance Ratio por arreglo, con el modelo bifacial. \\
\cd{irradiancia} & GHI, kt* e insolación, con control de calidad. \\
\cd{temperatura} & Temperatura por arreglo. \\
\cd{cobertura} & Qué datos hay: rango temporal y conteos. \\
\cd{catalogo} & Diccionario de variables. \\
\cd{tendencia} & Evolución de una variable en el tiempo. \\
\cd{graficar} & Serie lista para dibujar. \\
\bottomrule
\end{tabular}
\end{center}

El registro (\cd{tools/\_\_init\_\_.py}) es puro ensamblado: junta los esquemas que ve el
modelo y el mapa nombre\,$\to$\,función. Agregar una herramienta es crear su archivo e
importarlo ahí; no hay lógica que actualizar en dos lugares.

\subsection{Endpoints}

\begin{center}
\small
\begin{tabular}{@{}p{40mm}p{78mm}p{14mm}@{}}
\toprule
\textbf{Ruta} & \textbf{Para qué} & \textbf{Clave} \\
\midrule
\cd{GET /health}, \cd{/tools} & Ping y esquemas de las herramientas. & no \\
\cd{POST /tool/<nombre>} & Ejecutar una herramienta directo, sin LLM. & sí \\
\cd{POST /preguntar}, \cd{/chat} & Lazo completo del LLM, con traza. & sí \\
\cd{GET /datos/*} & Vistazo de sólo lectura a las relaciones permitidas. & sí \\
\bottomrule
\end{tabular}
\end{center}

Exponer cada herramienta como endpoint HTTP tiene un motivo concreto: permite verificar
un número sin gastar tokens, y permite que un orquestador externo la cablee como nodo de
acción en vez de depender del lazo conversacional.

\begin{aviso}[title=El borde de seguridad de \cd{/datos/*}]
Esos endpoints interpolan nombres de tabla dentro del SQL, porque un identificador no se
puede parametrizar. Por eso la \textbf{lista blanca} de \cd{datos.py} es el borde de
seguridad real: agregar una relación es agregarla ahí, nunca aceptar el nombre desde
afuera. Hay pruebas que fijan esa regla para que no se relaje por descuido.
\end{aviso}

\section{Consola de depuración}

Una web mínima en Next.js para probar los dos agentes en vivo, con datos reales. No es
un producto: es el instrumento con el que se verifica que un agente no está inventando.

\subsection{La regla de verificación}

Por cada pregunta, la consola dibuja la traza completa: cada turno del modelo con su
texto y las herramientas que decide llamar, cada ejecución real con su entrada, su
\textbf{salida cruda}, su tiempo en milisegundos y su error si lo hubo, y al final la
respuesta redactada con el consumo de tokens y el costo.

\begin{clave}[title=Regla de oro]
Todo número de la respuesta final tiene que aparecer en la salida de alguna herramienta.
Si aparece un número que no está en ninguna salida, el modelo lo alucinó.
\end{clave}

\subsection{Camino de un request}

\begin{center}
\small
\cd{Navegador} $\to$ \cd{/api/analizador/*} $\to$ \cd{:8010} $\to$ Supabase (RO) \\[2pt]
\cd{Navegador} $\to$ \cd{/api/pronostico/*} $\to$ \cd{:8000} $\to$ store (RO)
\end{center}

El navegador \textbf{nunca} habla directo con los servicios Python ni ve las API keys.
Todo pasa por manejadores de ruta del lado servidor, que reenvían la petición inyectando
la clave. Las claves viven sólo en la configuración del servidor.

Los endpoints que la consola necesita se agregaron a los propios agentes, no se
reimplementaron en Node. Es una decisión de fuente única de verdad: si el lazo del LLM se
reimplementara en el cliente, habría dos lazos que mantener sincronizados y la traza
dejaría de reflejar lo que realmente corre en producción.

\subsection{Control de acceso}

El \emph{middleware} decide si un request pasa; no sabe firmar cookies ni renderizar el
login. Las páginas redirigen al login; las rutas \cd{/api/*} responden 401 en JSON,
porque un \cd{fetch} no puede hacer nada útil con una redirección a HTML.

\begin{aviso}[title=Fallar cerrado]
Sin \cd{DEBUGGER\_PASSWORD} configurada, en producción la consola responde \textbf{503}
en lugar de abrirse. Fallar abierto fue exactamente lo que la dejó expuesta en Amplify
en su momento. En desarrollo sí deja pasar, para no estorbar el trabajo local.
\cd{DEBUGGER\_SESSION\_SECRET} firma la cookie: rotarlo cierra todas las sesiones sin
cambiarle la contraseña al equipo.
\end{aviso}

\subsection{Dos formas de levantarla}

\cd{./dev.sh} monta todo local: analizador, pronóstico y la web. \cd{./consola.sh} abre
un túnel SSH contra los agentes que \emph{ya corren} en la EC2 ---que siguen escuchando
sólo en el loopback del servidor, sin exponer nada--- elige puertos libres (8000, 8010 y
3000 suelen estar ocupados en una máquina de trabajo), sincroniza las URLs con los
puertos de esa corrida y cierra el túnel al salir.

\section{Despliegue y operación}

\begin{figure}[htbp]
\centering
\resizebox{0.96\textwidth}{!}{% Figura: despliegue real (desde 2026-08-25). Los servicios estan en loopback y lo
% unico publicado es nginx con TLS. VisioneFlow consume la API como un tercero mas.
\begin{tikzpicture}[node distance=7mm]

% --- Quien consume, desde internet ------------------------------------------
\node[compx,text width=34mm] (nav) {Navegador\\\nota{persona del equipo}};
\node[compx,text width=34mm,below=14mm of nav] (vf) {VisioneFlow\\\nota{nodo de agente}\\\nota{y disparo cada hora}};
\node[grupo,fit=(nav)(vf),label={[titgrupo]above:INTERNET}] (gnet) {};

% --- El servidor ------------------------------------------------------------
\node[comp,text width=30mm,right=26mm of nav,yshift=-9mm] (ngx)
  {nginx\\\nota{TLS, puertos 80 y 443}\\\nota{limita el acceso}};

\node[comp,font=\scriptsize,text width=44mm,right=30mm of ngx,yshift=16mm] (con)
  {Consola\\\nota{\cd{127.0.0.1:3001}}\\\nota{contraseña y cookie firmada}};
\node[comp,font=\scriptsize,text width=44mm,below=of con] (pre)
  {\cd{predictivo-predictivo-1}\\\nota{\cd{127.0.0.1:8000}}};
\node[comp,font=\scriptsize,text width=44mm,below=of pre] (his)
  {\cd{historico-historico-1}\\\nota{\cd{127.0.0.1:8010}}};
\node[agro,font=\scriptsize,text width=44mm,below=of his] (pg)
  {\cd{agrodash-pg}\\\nota{réplica del volcado}\\\nota{\cd{127.0.0.1:5433}}};

\node[comp,font=\scriptsize,text width=32mm,below=14mm of ngx] (tim)
  {temporizadores systemd\\\nota{\cd{predictivo-etl} 15 min}\\\nota{\cd{predictivo-refresh} 6 h}};

\node[grupo,fit=(ngx)(con)(pre)(his)(pg)(tim),
      label={[titgrupo]above:EC2 VisioneMetrics, apagada de 19:00 a 07:00}] (gsrv) {};

% --- Almacenamiento ---------------------------------------------------------
\node[almacen,font=\scriptsize,text width=34mm,right=14mm of his] (sb)
  {Supabase\\\nota{RLS activo}};

% --- Flujos -----------------------------------------------------------------
\draw[flecha] (nav.east) -- node[etiq,above,sloped]{HTTPS} (ngx.west);
\draw[flecha] (vf.east)  -- node[etiq,below,sloped]{HTTPS + clave} (ngx.west);

\draw[flecha] (ngx.east) -- node[etiq,pos=0.62,above,sloped]{\cd{/}} (con.west);
\draw[flecha] (ngx.east) -- node[etiq,pos=0.68,above]{\cd{/predictivo/}} (pre.west);
\draw[flecha] (ngx.east) -- node[etiq,pos=0.72,below,sloped]{\cd{/historico/}} (his.west);

\draw[flechad] (con.south) -- node[etiq,right]{loopback} (pre.north);
\draw[flecha] (pg.east) -- ++(6mm,0) |- node[etiq,pos=0.25,right]{fuente} (pre.east);
\draw[flecha] (tim.south) -- ++(0,-5mm) -| node[etiq,pos=0.25,below]{\cd{docker exec}} (pre.south);

\draw[flecha] (pre.east) -- ++(8mm,0) |- (sb.north);
\draw[flecha] (his.east) -- (sb.west);

\end{tikzpicture}
}
\caption{Qué corre dónde. Los dos agentes y la réplica escuchan sólo en el
\emph{loopback}; lo único publicado es nginx, con TLS, y detrás de él la consola con su
control de acceso y las dos rutas de servicio que consume VisioneFlow.}
\end{figure}

\subsection{Por qué la plataforma vive en su propio servidor}

Hasta el 2026-08-24 los dos microservicios corrían \emph{dentro} de la EC2 de VisioneFlow,
la plataforma de agentes de la empresa. Funcionaba, pero ataba dos proyectos que no
tienen nada que ver: el despliegue del pronóstico vivía en el repositorio de la
plataforma, cada despliegue del \emph{backend} hacía \cd{git reset --hard} y había que
acordarse de no tocar ciertos archivos, y una caída de cualquiera de los dos arrastraba
al otro.

Desde entonces la plataforma de datos tiene servidor propio, dominio propio y
certificado propio. VisioneFlow le sigue hablando, pero ahora como lo que es: una API
externa, con su clave, sobre HTTPS. Es la misma relación que tendría con cualquier
proveedor de terceros, y esa es exactamente la idea.

\subsection{Qué corre dónde}

\begin{center}
\small
\begin{tabular}{@{}p{34mm}p{50mm}p{56mm}@{}}
\toprule
\textbf{Pieza} & \textbf{Dónde} & \textbf{Cómo sobrevive a un reinicio} \\
\midrule
Agente Predictivo & \cd{predictivo-predictivo-1} \newline \nota{\cd{127.0.0.1:8000}} & \cd{restart: unless-stopped} \\
Agente Histórico & \cd{historico-historico-1} \newline \nota{\cd{127.0.0.1:8010}} & ídem \\
Réplica de AgroDash & \cd{agrodash-pg} \newline \nota{\cd{127.0.0.1:5433}} & ídem, sobre volumen persistente \\
Consola & proceso Node \newline \nota{\cd{127.0.0.1:3001}} & gestor de procesos con estado guardado \\
Borde HTTPS & nginx del sistema \newline \nota{puertos 80 y 443} & unidad systemd habilitada \\
Ingesta cada 15 min & \cd{predictivo-etl.timer} & unidad systemd habilitada \\
Refresco cada 6 h & \cd{predictivo-refresh.timer} & ídem \\
Almacenamiento & Supabase & servicio gestionado \\
\bottomrule
\end{tabular}
\end{center}

\begin{aviso}[title=El servidor no está encendido las 24 horas]
Un programador de tareas de AWS apaga la instancia a las 19:00 y la enciende a las 07:00,
de lunes a viernes. De viernes a lunes está apagada el fin de semana completo. Es una
decisión de costo, no un fallo, pero condiciona todo lo demás: cualquier consumidor
externo tiene que disparar dentro de esa ventana. La ingesta no pierde nada porque sus
temporizadores llevan \cd{Persistent=true} y \cd{OnBootSec}, así que al encender se pone
al día sola con lo que se perdió.
\end{aviso}

\subsection{Qué se publica, y qué no}

Sólo hay tres cosas alcanzables desde internet, todas por el mismo nginx con TLS:

\begin{center}
\small
\begin{tabular}{@{}p{40mm}p{40mm}p{58mm}@{}}
\toprule
\textbf{Ruta} & \textbf{Va a} & \textbf{Quién puede} \\
\midrule
\cd{/} & la consola & Quien tenga la contraseña. Sin sesión válida, las
  páginas redirigen y \cd{/api/*} responde 401. \\
\cd{/predictivo/...} & el Predictivo & Quien tenga la clave de servicio. No pasa por el
  control de la consola: un flujo automático no tiene sesión. \\
\cd{/historico/...} & el Histórico & ídem. \\
\bottomrule
\end{tabular}
\end{center}

Los nombres anteriores (\cd{/forecast/} y \cd{/analizador/}) siguen respondiendo como
alias, para que el cambio de servidor no obligara a reconfigurar rutas en el canvas de
la plataforma.

\subsection{Por qué hay una réplica de AgroDash}

El servidor de Cartago está caído, así que la fuente de la ingesta pasó a ser una
réplica restaurada de un volcado, montada en el propio servidor. El puerto queda atado al
loopback y ocupa unos 6\,GB. Volver a la fuente viva es cambiar una línea de
configuración: la URL de Cartago quedó comentada, no borrada.

Un detalle que cuesta tiempo si no se sabe: si el \cd{pg\_restore} de la máquina es más
nuevo que el servidor, falla con \cd{unrecognized configuration parameter "transaction\_timeout"}.
Por eso la restauración se hace con el \cd{pg\_restore} de \emph{adentro} del contenedor.

Al mudar el servidor, la réplica se movió por la red privada de la nube, sin salir a
internet, y se verificó contando las filas de las dos puntas: 21\,314\,662 lecturas en el
origen y las mismas 21\,314\,662 en el destino, con idéntica marca de tiempo máxima. Un
\cd{pg\_restore} que termina con código cero no es prueba de nada; el conteo sí.

\subsection{Salud del sistema}

Se mira sin entrar al servidor, por HTTP: \cd{/salud/ingesta} da la antigüedad de los
datos por variable (503 si están viejos) y \cd{/salud/panel} agrega errores recientes,
gasto del día y última predicción. La consola los muestra en su sección de salud.

\begin{aviso}[title=El 503 actual es correcto]
La fuente de San Carlos está congelada desde el 2026-07-23: el subsistema de cajas con
sufijo \cd{SC} y los nodos de suelo dejaron de reportar. \cd{/salud/ingesta} responde 503
porque está informando la antigüedad \emph{real} de los datos, no porque el agente falle.
Todo corre en modo ``sólo histórico''; cuando el equipo restaure la ingesta, el ETL
rellena solo por ser idempotente e incremental.
\end{aviso}

\subsection{Lección de observabilidad}

El ETL falló cada 15 minutos durante nueve días \textbf{sin dejar rastro} en
\cd{agente\_log}. La causa: la conexión a la fuente estaba \emph{fuera} del bloque
\cd{try}, así que la excepción se propagaba antes de que hubiera con qué registrarla.
La corrección tiene tres partes, y las tres importan: conectar primero al store (que es
donde se escribe el log), meter la conexión a la fuente dentro del \cd{try} con su propio
evento de fallo, y \textbf{volver a lanzar} la excepción para que systemd siga marcando
la unidad como fallida. Hay una prueba de regresión que fija ese comportamiento.

\subsection{Diagnóstico rápido}

\begin{center}
\small
\begin{tabular}{@{}p{56mm}p{84mm}@{}}
\toprule
\textbf{Síntoma} & \textbf{Dónde mirar} \\
\midrule
La ingesta no trae datos & \cd{systemctl status predictivo-etl.service} y la tabla
  \cd{agente\_log} filtrando por nivel de error. \\
El pronóstico responde datos viejos & \cd{GET /salud/ingesta}: la fuente está congelada,
  no es un fallo del agente. \\
Nada responde, de noche o en fin de semana & El servidor está apagado por programación.
  No es una caída. \\
Todo responde 429 & Tope diario de gasto, límite de ritmo del servicio, o el freno de
  nginx si es el acceso a la consola. \\
La consola responde 503 & Falta \cd{DEBUGGER\_PASSWORD} en el entorno. Falla cerrada a
  propósito. \\
Un servicio da 502 & El contenedor está caído; nginx sigue en pie. \\
Supabase rechaza escrituras & El plan gratuito son 500\,MB y está cerca del límite. \\
\bottomrule
\end{tabular}
\end{center}

\subsection{Pruebas}

Ninguna suite necesita credenciales ni red hacia Supabase; si una empieza a pedirlas,
dejó de ser unitaria. Corren en CI en cada \emph{push} y cada PR.

\begin{center}
\small
\begin{tabular}{@{}p{40mm}p{26mm}p{68mm}@{}}
\toprule
\textbf{Componente} & \textbf{Pruebas} & \textbf{Qué cubren} \\
\midrule
Agente Predictivo & 253 & Horizonte, física, herramienta, ETL, límites, salud, claves,
  riesgo, climatología, observabilidad, API. \\
Agente Histórico & 61 & API, costos, lista blanca de datos, control de calidad, cielo,
  contexto. \\
Consola & 3 guiones de humo & Construcción, control de acceso y formato de respuestas. \\
ETL de CSV & ninguna & Sin suite; la verificación es el \emph{notebook} de EDA. \\
\bottomrule
\end{tabular}
\end{center}

\section{Seguridad}

\begin{itemize}[leftmargin=*,itemsep=4pt]
  \item \textbf{AgroDash es sólo lectura, y se fuerza.} Toda conexión fija
    \cd{SET SESSION CHARACTERISTICS AS TRANSACTION READ ONLY} y sólo se ejecutan
    \cd{SELECT}. No se confía en que el rol esté bien configurado del otro lado.

  \item \textbf{RLS en modo cierre} en las tablas públicas de Supabase, deliberadamente
    \emph{sin políticas}: sólo los roles de servicio acceden y la API REST pública queda
    bloqueada. Se verificó que no rompe al forecaster, que corre con un rol que la
    salta.

  \item \textbf{Vistas con \cd{security\_invoker}}, para que ejecuten con los permisos de
    quien consulta y no con los del dueño.

  \item \textbf{Claves sólo por entorno}, nunca en el código ni en el repositorio.
    Ningún secreto está versionado.

  \item \textbf{Comparación en tiempo constante} de las API keys
    (\cd{secrets.compare\_digest}), para no filtrar la clave por diferencias de tiempo de
    respuesta.

  \item \textbf{Clave exigida sólo si está configurada.} En local los servicios corren
    abiertos para no estorbar; en producción la variable está definida y el borde exige
    el encabezado. La excepción son \cd{/health} y \cd{/salud/ingesta}, abiertas a
    propósito para monitoreo automático: no exponen datos de la serie.

  \item \textbf{Los servicios no se publican; se publica el borde.} Los dos agentes y la
    réplica escuchan en \cd{127.0.0.1} y nada más. Lo único alcanzable desde internet es
    nginx, con TLS. Hasta la mudanza los servicios escuchaban en \cd{0.0.0.0} y lo único
    que los tapaba era el grupo de seguridad de la nube: cualquier máquina de la misma
    red privada los alcanzaba. Atarlos al loopback quita esa clase entera de error.

  \item \textbf{Freno de fuerza bruta en el borde.} El límite de intentos de la propia
    consola vive en memoria del proceso, así que un reinicio lo borra y en un despliegue
    con varias instancias no se comparte. La defensa que no depende de eso es doble: cada
    intento cuesta un derivado de clave deliberadamente caro (unos 400\,ms de CPU), y
    nginx limita la ruta de acceso a seis intentos por minuto. Verificado, no sólo
    configurado: diez intentos seguidos dan seis 401 y cuatro 429.

  \item \textbf{Falla cerrada.} Sin contraseña configurada, la consola en producción no
    se abre: responde 503 y lo dice. Un despliegue anterior quedó expuesto justamente por
    hacer lo contrario, abrirse cuando faltaba la variable.

  \item \textbf{Errores internos genéricos.} Un fallo interno responde el 500 estándar;
    el detalle queda en el log del servidor, nunca en la respuesta.
\end{itemize}

\subsection{Secretos necesarios}

\begin{center}
\small
\begin{tabular}{@{}p{44mm}p{96mm}@{}}
\toprule
\textbf{Variable} & \textbf{Para qué} \\
\midrule
\cd{DATABASE\_URL} & Conexión a Supabase (usar el \emph{Session pooler}; la conexión
  directa es sólo IPv6). \\
\cd{STORE\_URL} & El store del agente, separado de la fuente a propósito. \\
\cd{ANTHROPIC\_API\_KEY} & Los dos agentes. \\
\cd{FORECAST\_API\_KEY} & Proteger el servicio de pronóstico. \\
\cd{HISTORICO\_API\_KEY} & Proteger el servicio del Histórico. \\
\cd{DEBUGGER\_PASSWORD} & Entrar a la consola. \\
\cd{DEBUGGER\_SESSION\_SECRET} & Firmar la cookie de sesión. \\
\bottomrule
\end{tabular}
\end{center}

\subsection{Pendiente de seguridad}

Rotar la clave débil del rol de sólo lectura que se usó para probar el acceso a la base
viva de Cartago por la malla Tailscale.

\section{Estado actual}

\subsection{Volúmenes actuales}

\begin{center}
\small
\begin{tabular}{@{}p{54mm}p{28mm}p{56mm}@{}}
\toprule
\textbf{Relación} & \textbf{Filas} & \textbf{Nota} \\
\midrule
\cd{monitoreo\_sc\_electrico} & 36\,469 & 285 CSV, ventanas de 5\,min. \\
\cd{radiacion\_sc\_15s} & 94\,868 & Ventanas de 15\,s. \\
\cd{radiacion\_sc\_clearsky} & 94\,868 & Un cielo despejado por instante. \\
\cd{radiacion\_sc\_poa} & 56\,450 & Instantes con POA por arreglo. \\
\cd{lecturas\_ambientales\_sc} & $\approx$\,812\,000 & Ingesta de AgroDash, formato largo. \\
\bottomrule
\end{tabular}
\end{center}

La verificación cruzada confirmó que el crudo sigue siendo crudo (temperaturas en 85 y el
\emph{offset} presentes en la tabla base) y que las vistas hacen su trabajo (85 pasa a
NULL, la potencia máxima cae a un valor físico, la irradiancia negativa queda en cero,
kt* con percentil 95 en 1{,}01 y 99{,}4\% de filas aprobando el control de calidad).

\subsection{Decisiones que conviene no revisitar}

\begin{itemize}[leftmargin=*,itemsep=3pt]
  \item El crudo se guarda; la corrección vive en vistas. Cualquier propuesta de
    ``limpiar al insertar'' ya se descartó explícitamente.
  \item No se generan datos sintéticos para los huecos largos de 126 y 71 días. Como
    referencia paralela se usa NASA POWER.
  \item Los duplicados exactos por md5 se eliminan; los fragmentos numerados requieren
    decisión humana.
  \item Las bases de las dos regiones no se fusionan a nivel de almacenamiento, aunque un
    agente lea ambas.
\end{itemize}

\vspace{1em}
\begin{clave}[title=Dónde seguir]
El sistema de memoria del proyecto (\cd{docs/memoria/INDEX.md}) es la fuente viva: un
tema por archivo, con la evidencia contada y la fecha de cada decisión. Este documento
resume la arquitectura; ese índice explica \emph{por qué} cada pieza quedó así.
\end{clave}


\end{document}
